\noindent
{\Large \textbf{\section*{Заключение}}} \par
\addcontentsline{toc}{section}{Заключение} \par

В ходе выполнения курсовой работы была выбрана предметная область «Организация выставочной деятельности музея». 
При ее анализе было задействовано 4 источника, указанные в списке использованных источников.
По результатам данного анализа были сформированы 3 цели функционирования базы данных:
ведение централизованного реестра музейных предметов, ведение реестра проводимых в музее выставок, хранение планов размещения выставок в залах музея.
Также было выделено 4 роли и 10 сущностей. На основании выделенных сущностей была построена ER-диаграмма в нотации Чена и схема связи объектов при помощи онлайн-ресурса draw.io.

На этапе проектирования было выделено 20 таблиц. Для атрибутов были выбраны типы SERIAL, INT, VARCHAR, DECIMAL, при этом выбор каждого типа был обоснован, опираясь на примеры и источники. 
Далее были построены схемы базы данных на русском и английском языках, а также схема заполнения уровней БД в draw.io.

На этапе программирования работа велась на ноутбуке Macbook M4 pro с 16гб оперативной памяти. В качестве IDE использовалась Visual Studio Code, для работы с БД использовалось СУБД PostgreSQL.
Далее при помощи SQL-запросов было создано 20 таблиц, их заполнение было реализовано при помощи языка Python, всего в таблицах 2 646 202 записей. После этого было реализовано 9 запросов на языке SQL.
Для каждого запроса была построена реляционная формула и план выполнения.
Для результатов запросов №3 и №5 дополнительно были построены 2-d гистограммы, а для запроса №8 - 3-d гистограмма.


\newpage